% Options for packages loaded elsewhere
\PassOptionsToPackage{unicode}{hyperref}
\PassOptionsToPackage{hyphens}{url}
%
\documentclass[
  11pt,
]{article}
\usepackage{amsmath,amssymb}
\usepackage{iftex}
\ifPDFTeX
  \usepackage[T1]{fontenc}
  \usepackage[utf8]{inputenc}
  \usepackage{textcomp} % provide euro and other symbols
\else % if luatex or xetex
  \usepackage{unicode-math} % this also loads fontspec
  \defaultfontfeatures{Scale=MatchLowercase}
  \defaultfontfeatures[\rmfamily]{Ligatures=TeX,Scale=1}
\fi
\usepackage{lmodern}
\ifPDFTeX\else
  % xetex/luatex font selection
\fi
% Use upquote if available, for straight quotes in verbatim environments
\IfFileExists{upquote.sty}{\usepackage{upquote}}{}
\IfFileExists{microtype.sty}{% use microtype if available
  \usepackage[]{microtype}
  \UseMicrotypeSet[protrusion]{basicmath} % disable protrusion for tt fonts
}{}
\makeatletter
\@ifundefined{KOMAClassName}{% if non-KOMA class
  \IfFileExists{parskip.sty}{%
    \usepackage{parskip}
  }{% else
    \setlength{\parindent}{0pt}
    \setlength{\parskip}{6pt plus 2pt minus 1pt}}
}{% if KOMA class
  \KOMAoptions{parskip=half}}
\makeatother
\usepackage{xcolor}
\usepackage[margin=2.2cm]{geometry}
\usepackage{longtable,booktabs,array}
\usepackage{calc} % for calculating minipage widths
% Correct order of tables after \paragraph or \subparagraph
\usepackage{etoolbox}
\makeatletter
\patchcmd\longtable{\par}{\if@noskipsec\mbox{}\fi\par}{}{}
\makeatother
% Allow footnotes in longtable head/foot
\IfFileExists{footnotehyper.sty}{\usepackage{footnotehyper}}{\usepackage{footnote}}
\makesavenoteenv{longtable}
\setlength{\emergencystretch}{3em} % prevent overfull lines
\providecommand{\tightlist}{%
  \setlength{\itemsep}{0pt}\setlength{\parskip}{0pt}}
\setcounter{secnumdepth}{-\maxdimen} % remove section numbering
\ifLuaTeX
  \usepackage{selnolig}  % disable illegal ligatures
\fi
\IfFileExists{bookmark.sty}{\usepackage{bookmark}}{\usepackage{hyperref}}
\IfFileExists{xurl.sty}{\usepackage{xurl}}{} % add URL line breaks if available
\urlstyle{same}
\hypersetup{
  pdftitle={Worksheet 4A: Derivatives and Integrals --- Inverse Operations},
  hidelinks,
  pdfcreator={LaTeX via pandoc}}

\title{Worksheet 4A: Derivatives and Integrals --- Inverse Operations}
\author{}
\date{}

\begin{document}
\maketitle

\begin{quote}
\textbf{What This Worksheet Is About}

Derivatives and integrals are \textbf{inverses} of each other, just like
multiplication and division, or squaring and square roots.

\begin{itemize}
\tightlist
\item
  \textbf{Derivative:} Given a function, find its rate of change
\item
  \textbf{Integral:} Given a rate of change, find the original function
\end{itemize}

But there's a twist: When we ``undo'' a derivative, we don't get back
exactly one function --- we get a whole \textbf{family} of functions!
\end{quote}

\begin{center}\rule{0.5\linewidth}{0.5pt}\end{center}

\hypertarget{part-a-the-power-rule-both-directions-20-problems}{%
\subsection{Part A: The Power Rule --- Both Directions (20
problems)}\label{part-a-the-power-rule-both-directions-20-problems}}

\hypertarget{differentiation-bring-down-the-power-reduce-by-1}{%
\subsubsection{Differentiation: Bring down the power, reduce by
1}\label{differentiation-bring-down-the-power-reduce-by-1}}

\[\frac{d}{dx}x^n = n \cdot x^{n-1}\]

\hypertarget{integration-increase-the-power-by-1-divide-by-new-power}{%
\subsubsection{Integration: Increase the power by 1, divide by new
power}\label{integration-increase-the-power-by-1-divide-by-new-power}}

\[\int x^n \, dx = \frac{x^{n+1}}{n+1} + C \quad \text{(for } n \neq -1\text{)}\]

\textbf{Notice they're inverses!} - Derivative: exponent comes
\textbf{down} and \textbf{decreases} - Integral: exponent goes
\textbf{up} and we \textbf{divide} by it

\begin{center}\rule{0.5\linewidth}{0.5pt}\end{center}

\textbf{Find both the derivative and the integral of each function:}

\textbf{1.} \(f(x) = x^4\)

\begin{itemize}
\tightlist
\item
  \(f'(x) ={}\) \_\_\_\_\_\_\_\_\_\_
\item
  \(\int x^4 \, dx ={}\) \_\_\_\_\_\_\_\_\_\_
\end{itemize}

\textbf{2.} \(g(t) = t^5\)

\begin{itemize}
\tightlist
\item
  \(g'(t) ={}\) \_\_\_\_\_\_\_\_\_\_
\item
  \(\int t^5 \, dt ={}\) \_\_\_\_\_\_\_\_\_\_
\end{itemize}

\textbf{3.} \(h(p) = p^3\)

\begin{itemize}
\tightlist
\item
  \(h'(p) ={}\) \_\_\_\_\_\_\_\_\_\_
\item
  \(\int p^3 \, dp ={}\) \_\_\_\_\_\_\_\_\_\_
\end{itemize}

\textbf{4.} \(y = x^7\)

\textbf{5.} \(f(m) = m^2\)

\textbf{6.} \(g(r) = r^{10}\)

\textbf{7.} \(h(k) = k^1 = k\)

\textbf{8.} \(y = x^{100}\)

\begin{center}\rule{0.5\linewidth}{0.5pt}\end{center}

\textbf{Negative and Fractional Exponents:}

\textbf{9.} \(f(x) = x^{-2}\)

\begin{itemize}
\tightlist
\item
  \(f'(x) ={}\) \_\_\_\_\_\_\_\_\_\_
\item
  \(\int x^{-2} \, dx ={}\) \_\_\_\_\_\_\_\_\_\_
\end{itemize}

\textbf{10.} \(g(t) = t^{-3}\)

\textbf{11.} \(h(x) = x^{1/2} = \sqrt{x}\)

\begin{itemize}
\tightlist
\item
  \(h'(x) ={}\) \_\_\_\_\_\_\_\_\_\_
\item
  \(\int x^{1/2} \, dx ={}\) \_\_\_\_\_\_\_\_\_\_
\end{itemize}

\textbf{12.} \(y = x^{1/3}\)

\textbf{13.} \(f(p) = p^{-1/2}\)

\textbf{14.} \(g(n) = n^{3/2}\)

\begin{center}\rule{0.5\linewidth}{0.5pt}\end{center}

\textbf{The Special Case: What happens when \(n = -1\)?}

\textbf{15.} Try to compute \(\int x^{-1} \, dx\) using the formula
\(\frac{x^{n+1}}{n+1}\).

What goes wrong?

\textbf{16.} The integral of \(x^{-1} = \frac{1}{x}\) is actually
\(\ln|x| + C\). We'll explore why in Worksheet 4D!

\begin{center}\rule{0.5\linewidth}{0.5pt}\end{center}

\textbf{Checking Your Work --- Derivatives and Integrals Undo Each
Other:}

\textbf{17.} If \(\int x^3 \, dx = \frac{x^4}{4} + C\), verify by
differentiating: \(\frac{d}{dx}\left(\frac{x^4}{4} + C\right) ={}\)
\_\_\_\_\_

\textbf{18.} If \(\frac{d}{dx}(x^5) = 5x^4\), verify by integrating:
\(\int 5x^4 \, dx ={}\) \_\_\_\_\_

\textbf{19.} Differentiate, then integrate: Start with \(f(x) = x^6\) -
\(f'(x) ={}\) \_\_\_\_\_ - \(\int f'(x) \, dx ={}\) \_\_\_\_\_ - Do you
get back \(x^6\)? (Plus what?)

\textbf{20.} Integrate, then differentiate: Start with \(g(x) = x^4\) -
\(\int g(x) \, dx ={}\) \_\_\_\_\_ -
\(\frac{d}{dx}\left[\int g(x) \, dx\right] ={}\) \_\_\_\_\_ - Do you get
back \(x^4\)?

\begin{center}\rule{0.5\linewidth}{0.5pt}\end{center}

\hypertarget{part-b-why-c-the-family-of-antiderivatives-14-problems}{%
\subsection{Part B: Why ``+C''? --- The Family of Antiderivatives (14
problems)}\label{part-b-why-c-the-family-of-antiderivatives-14-problems}}

Consider these three functions: - \(f(x) = x^2\) - \(g(x) = x^2 + 5\) -
\(h(x) = x^2 - 17\)

\textbf{21.} Find the derivative of each: - \(f'(x) ={}\) \_\_\_\_\_ -
\(g'(x) ={}\) \_\_\_\_\_ - \(h'(x) ={}\) \_\_\_\_\_

\textbf{22.} What do you notice? Why does this happen?

\textbf{23.} The functions \(x^2\), \(x^2 + 5\), and \(x^2 - 17\) are
all different, but they have the \_\_\_\_\_\_\_\_\_\_ derivative.

\textbf{24.} Graphically: Sketch \(y = x^2\), \(y = x^2 + 3\), and
\(y = x^2 - 2\) on the same axes.

\begin{itemize}
\tightlist
\item
  How are these parabolas related?
\item
  At any given x-value, do they have the same slope?
\end{itemize}

\begin{center}\rule{0.5\linewidth}{0.5pt}\end{center}

\textbf{The Key Insight:}

\begin{quote}
When we integrate \(2x\), we get \(x^2 + C\) where \(C\) is \textbf{any
constant}.

This represents a whole \textbf{family} of functions --- infinitely many
parabolas, all with the same shape but shifted up or down.
\end{quote}

\begin{center}\rule{0.5\linewidth}{0.5pt}\end{center}

\textbf{25.} Write three different functions whose derivative is
\(3x^2\).

\textbf{26.} Write three different functions whose derivative is
\(5t^4\).

\textbf{27.} Write the general antiderivative (with +C) of \(4p^3\).

\textbf{28.} If \(F'(x) = 6x^5\), what is \(F(x)\)?

\begin{center}\rule{0.5\linewidth}{0.5pt}\end{center}

\textbf{Finding a Specific Function --- Initial Conditions:}

If we know \(F'(x) = 2x\) AND \(F(0) = 7\), we can find the specific
function:

\(F(x) = x^2 + C\)

\(F(0) = 0^2 + C = 7\), so \(C = 7\)

Therefore \(F(x) = x^2 + 7\)

\begin{center}\rule{0.5\linewidth}{0.5pt}\end{center}

\textbf{29.} Find \(F(x)\) if \(F'(x) = 3x^2\) and \(F(0) = 5\).

\textbf{30.} Find \(G(t)\) if \(G'(t) = 4t^3\) and \(G(1) = 10\).

\textbf{31.} Find \(H(p)\) if \(H'(p) = 2p + 1\) and \(H(2) = 0\).

\textbf{32.} A ball's velocity is \(v(t) = 20 - 10t\) m/s. Its position
is \(s(t) = \int v(t) \, dt\).

\begin{itemize}
\item
  \begin{enumerate}
  \def\labelenumi{(\alph{enumi})}
  \tightlist
  \item
    Find \(s(t)\) with the constant \(C\).
  \end{enumerate}
\item
  \begin{enumerate}
  \def\labelenumi{(\alph{enumi})}
  \setcounter{enumi}{1}
  \tightlist
  \item
    If the ball starts at height \(s(0) = 5\) meters, find \(C\).
  \end{enumerate}
\item
  \begin{enumerate}
  \def\labelenumi{(\alph{enumi})}
  \setcounter{enumi}{2}
  \tightlist
  \item
    Write the complete position function.
  \end{enumerate}
\end{itemize}

\textbf{33.} An object's acceleration is \(a(t) = 6\) m/s². Starting
from rest (\(v(0) = 0\)):

\begin{itemize}
\item
  \begin{enumerate}
  \def\labelenumi{(\alph{enumi})}
  \tightlist
  \item
    Find velocity: \(v(t) = \int a(t) \, dt\)
  \end{enumerate}
\item
  \begin{enumerate}
  \def\labelenumi{(\alph{enumi})}
  \setcounter{enumi}{1}
  \tightlist
  \item
    If position \(s(0) = 10\), find \(s(t) = \int v(t) \, dt\).
  \end{enumerate}
\end{itemize}

\textbf{34.} Explain in your own words why integration has ``+C'' but
differentiation doesn't.

\begin{center}\rule{0.5\linewidth}{0.5pt}\end{center}

\hypertarget{part-c-constants-and-sums-both-directions-16-problems}{%
\subsection{Part C: Constants and Sums --- Both Directions (16
problems)}\label{part-c-constants-and-sums-both-directions-16-problems}}

\hypertarget{derivative-rules}{%
\subsubsection{Derivative Rules:}\label{derivative-rules}}

\begin{itemize}
\tightlist
\item
  \(\frac{d}{dx}(c) = 0\) (constant rule)
\item
  \(\frac{d}{dx}[cf(x)] = c \cdot f'(x)\) (constant multiple)
\item
  \(\frac{d}{dx}[f(x) + g(x)] = f'(x) + g'(x)\) (sum rule)
\end{itemize}

\hypertarget{integral-rules-the-reverses}{%
\subsubsection{Integral Rules (the
reverses):}\label{integral-rules-the-reverses}}

\begin{itemize}
\tightlist
\item
  \(\int 0 \, dx = C\) (integrating zero gives a constant)
\item
  \(\int c \cdot f(x) \, dx = c \cdot \int f(x) \, dx\) (constant
  multiple)
\item
  \(\int [f(x) + g(x)] \, dx = \int f(x) \, dx + \int g(x) \, dx\) (sum
  rule)
\end{itemize}

\textbf{Also:} \(\int c \, dx = cx + C\) (integrating a constant)

\begin{center}\rule{0.5\linewidth}{0.5pt}\end{center}

\textbf{Find both the derivative and the integral:}

\textbf{35.} \(f(x) = 5x^3\)

\begin{itemize}
\tightlist
\item
  \(f'(x) ={}\) \_\_\_\_\_\_\_\_\_\_
\item
  \(\int 5x^3 \, dx ={}\) \_\_\_\_\_\_\_\_\_\_
\end{itemize}

\textbf{36.} \(g(t) = -2t^4\)

\textbf{37.} \(h(p) = \frac{1}{2}p^6\)

\textbf{38.} \(y = 7\) (a constant function)

\begin{itemize}
\tightlist
\item
  \(\frac{dy}{dx} ={}\) \_\_\_\_\_\_\_\_\_\_
\item
  \(\int 7 \, dx ={}\) \_\_\_\_\_\_\_\_\_\_
\end{itemize}

\textbf{39.} \(f(m) = -3\)

\begin{center}\rule{0.5\linewidth}{0.5pt}\end{center}

\textbf{Sums and Differences:}

\textbf{40.} \(f(x) = x^3 + x^2\)

\begin{itemize}
\tightlist
\item
  \(f'(x) ={}\) \_\_\_\_\_\_\_\_\_\_
\item
  \(\int (x^3 + x^2) \, dx ={}\) \_\_\_\_\_\_\_\_\_\_
\end{itemize}

\textbf{41.} \(g(t) = t^4 - t^2\)

\textbf{42.} \(h(r) = 3r^2 + 2r - 5\)

\textbf{43.} \(y = 4x^3 - 6x^2 + 2x - 1\)

\textbf{44.} \(f(k) = k^5 - 3k^3 + 2k\)

\textbf{45.} \(g(n) = 2n^4 + n^3 - 5n^2 + 4n - 7\)

\begin{center}\rule{0.5\linewidth}{0.5pt}\end{center}

\textbf{Verify by reversing:}

\textbf{46.} Check problem 40: Differentiate your integral. Do you get
back \(x^3 + x^2\)?

\textbf{47.} Check problem 43: Integrate your derivative. Do you get
back \(4x^3 - 6x^2 + 2x - 1\) (plus C)?

\begin{center}\rule{0.5\linewidth}{0.5pt}\end{center}

\textbf{Rewriting Before Integrating:}

\textbf{48.} Find \(\int \frac{3}{x^2} \, dx\)

\emph{Hint: Rewrite as \(3x^{-2}\) first}

\textbf{49.} Find \(\int \frac{5}{\sqrt{t}} \, dt\)

\textbf{50.} Find \(\int \left(x + \frac{1}{x^2}\right) dx\)

\begin{center}\rule{0.5\linewidth}{0.5pt}\end{center}

\hypertarget{part-d-polynomials-complete-practice-16-problems}{%
\subsection{Part D: Polynomials --- Complete Practice (16
problems)}\label{part-d-polynomials-complete-practice-16-problems}}

\textbf{Find the derivative:}

\textbf{51.} \(f(x) = 2x^5 - 3x^4 + x^3 - 7x + 2\)

\textbf{52.} \(g(t) = \frac{t^4}{4} - \frac{t^3}{3} + \frac{t^2}{2}\)

\textbf{53.} \(h(p) = (p + 2)(p - 3)\) \emph{(expand first)}

\textbf{54.} \(y = (x - 1)^2\)

\begin{center}\rule{0.5\linewidth}{0.5pt}\end{center}

\textbf{Find the integral:}

\textbf{55.} \(\int (3x^2 - 4x + 5) \, dx\)

\textbf{56.} \(\int (t^4 - 2t^3 + t - 3) \, dt\)

\textbf{57.} \(\int (6m^5 - 4m^3 + 2m) \, dm\)

\textbf{58.} \(\int (u^2 + 1)^2 \, du\) \emph{(expand first!)}

\begin{center}\rule{0.5\linewidth}{0.5pt}\end{center}

\textbf{Mixed --- Derivative or Integral as Specified:}

\textbf{59.} \(\frac{d}{dx}(x^4 - 3x^2 + 2)\)

\textbf{60.} \(\int (4r^3 - 6r) \, dr\)

\textbf{61.} \(\frac{d}{dt}(t^6 + 2t^4 - t)\)

\textbf{62.} \(\int (5k^4 - 3k^2 + 1) \, dk\)

\begin{center}\rule{0.5\linewidth}{0.5pt}\end{center}

\textbf{Round-Trip Problems:}

\textbf{63.} Start with \(f(x) = x^5\). - (a) Differentiate:
\(f'(x) ={}\) - (b) Integrate your answer: \(\int f'(x) \, dx ={}\) -
(c) Is this \(f(x)\)? Explain.

\textbf{64.} Start with \(g(t) = t^3 - 4t\). - (a) Integrate:
\(\int g(t) \, dt ={}\) - (b) Differentiate your answer. - (c) Is this
\(g(t)\)? Explain.

\textbf{65.} Why does integrating then differentiating give you back
exactly what you started with, but differentiating then integrating
gives you the original plus a constant?

\textbf{66.} Complete this analogy:

Squaring and square roots are inverses: \((\sqrt{x})^2 = x\) and
\(\sqrt{x^2} = |x|\) (almost x, but not quite for negatives)

Similarly: \(\frac{d}{dx}\left[\int f(x) \, dx\right] ={}\) \_\_\_\_\_
and \(\int \frac{d}{dx}[f(x)] \, dx ={}\) \_\_\_\_\_

\begin{center}\rule{0.5\linewidth}{0.5pt}\end{center}

\hypertarget{part-e-building-intuition-whats-really-happening-10-problems}{%
\subsection{Part E: Building Intuition --- What's Really Happening (10
problems)}\label{part-e-building-intuition-whats-really-happening-10-problems}}

\textbf{67.} The derivative of position is velocity. The integral of
velocity is \_\_\_\_\_\_\_\_\_\_.

\textbf{68.} The derivative of velocity is acceleration. The integral of
acceleration is \_\_\_\_\_\_\_\_\_\_.

\textbf{69.} If you know an object's acceleration at all times, can you
determine its exact position? What additional information do you need?

\textbf{70.} Match the physical quantities:

\begin{longtable}[]{@{}lll@{}}
\toprule\noalign{}
Function & Derivative & Integral \\
\midrule\noalign{}
\endhead
\bottomrule\noalign{}
\endlastfoot
Position & & \\
Velocity & & \\
Acceleration & & \\
\end{longtable}

\begin{center}\rule{0.5\linewidth}{0.5pt}\end{center}

\textbf{71.} A car's velocity is \(v(t) = 3t^2\) m/s.

\begin{itemize}
\item
  \begin{enumerate}
  \def\labelenumi{(\alph{enumi})}
  \tightlist
  \item
    What is the acceleration \(a(t) = v'(t)\)?
  \end{enumerate}
\item
  \begin{enumerate}
  \def\labelenumi{(\alph{enumi})}
  \setcounter{enumi}{1}
  \tightlist
  \item
    What is the position \(s(t) = \int v(t) \, dt\)?
  \end{enumerate}
\item
  \begin{enumerate}
  \def\labelenumi{(\alph{enumi})}
  \setcounter{enumi}{2}
  \tightlist
  \item
    If \(s(0) = 0\), what is the specific position function?
  \end{enumerate}
\item
  \begin{enumerate}
  \def\labelenumi{(\alph{enumi})}
  \setcounter{enumi}{3}
  \tightlist
  \item
    How far has the car traveled from \(t = 0\) to \(t = 2\)?
  \end{enumerate}
\end{itemize}

\textbf{72.} An object's position is \(s(t) = t^3 - 6t^2 + 9t\).

\begin{itemize}
\item
  \begin{enumerate}
  \def\labelenumi{(\alph{enumi})}
  \tightlist
  \item
    Find velocity \(v(t) = s'(t)\).
  \end{enumerate}
\item
  \begin{enumerate}
  \def\labelenumi{(\alph{enumi})}
  \setcounter{enumi}{1}
  \tightlist
  \item
    Find acceleration \(a(t) = v'(t)\).
  \end{enumerate}
\item
  \begin{enumerate}
  \def\labelenumi{(\alph{enumi})}
  \setcounter{enumi}{2}
  \tightlist
  \item
    When is velocity zero?
  \end{enumerate}
\item
  \begin{enumerate}
  \def\labelenumi{(\alph{enumi})}
  \setcounter{enumi}{3}
  \tightlist
  \item
    When is acceleration zero?
  \end{enumerate}
\end{itemize}

\begin{center}\rule{0.5\linewidth}{0.5pt}\end{center}

\textbf{73.} Revenue \(R(q)\) depends on quantity sold \(q\).

\begin{itemize}
\tightlist
\item
  \textbf{Marginal revenue} is \(R'(q)\) --- the extra revenue from one
  more unit.
\item
  If you know marginal revenue, \(\int R'(q) \, dq\) gives you
  \_\_\_\_\_\_\_\_\_\_.
\end{itemize}

\textbf{74.} Cost \(C(q)\) depends on quantity produced.

\begin{itemize}
\tightlist
\item
  \textbf{Marginal cost} is \(C'(q)\).
\item
  If marginal cost is \(MC(q) = 6q + 10\), find the cost function
  \(C(q)\).
\item
  If fixed costs are \(C(0) = 500\), find the specific cost function.
\end{itemize}

\begin{center}\rule{0.5\linewidth}{0.5pt}\end{center}

\textbf{75.} Population growth rate is \(P'(t) = 1000 + 50t\) people per
year.

\begin{itemize}
\item
  \begin{enumerate}
  \def\labelenumi{(\alph{enumi})}
  \tightlist
  \item
    Find \(P(t)\), the population function.
  \end{enumerate}
\item
  \begin{enumerate}
  \def\labelenumi{(\alph{enumi})}
  \setcounter{enumi}{1}
  \tightlist
  \item
    If the population at \(t = 0\) is 50,000, find the specific
    function.
  \end{enumerate}
\item
  \begin{enumerate}
  \def\labelenumi{(\alph{enumi})}
  \setcounter{enumi}{2}
  \tightlist
  \item
    What is the population after 10 years?
  \end{enumerate}
\end{itemize}

\textbf{76.} Water flows into a tank at rate \(r(t) = 100 - 2t\) liters
per minute.

\begin{itemize}
\item
  \begin{enumerate}
  \def\labelenumi{(\alph{enumi})}
  \tightlist
  \item
    The amount of water \(W(t) = \int r(t) \, dt\). Find it.
  \end{enumerate}
\item
  \begin{enumerate}
  \def\labelenumi{(\alph{enumi})}
  \setcounter{enumi}{1}
  \tightlist
  \item
    If the tank starts empty, \(W(0) = 0\). Find the specific function.
  \end{enumerate}
\item
  \begin{enumerate}
  \def\labelenumi{(\alph{enumi})}
  \setcounter{enumi}{2}
  \tightlist
  \item
    How much water is in the tank after 10 minutes?
  \end{enumerate}
\item
  \begin{enumerate}
  \def\labelenumi{(\alph{enumi})}
  \setcounter{enumi}{3}
  \tightlist
  \item
    When does the tank stop filling (when \(r(t) = 0\))?
  \end{enumerate}
\end{itemize}

\begin{center}\rule{0.5\linewidth}{0.5pt}\end{center}

\hypertarget{part-f-the-fundamental-theorem-preview-6-problems}{%
\subsection{Part F: The Fundamental Theorem Preview (6
problems)}\label{part-f-the-fundamental-theorem-preview-6-problems}}

The \textbf{Fundamental Theorem of Calculus} connects derivatives and
integrals precisely:

\[\int_a^b f(x) \, dx = F(b) - F(a)\]

where \(F'(x) = f(x)\).

This says: To find the integral from \(a\) to \(b\), find ANY
antiderivative \(F\), then compute \(F(b) - F(a)\).

(The \(+C\) cancels out!)

\begin{center}\rule{0.5\linewidth}{0.5pt}\end{center}

\textbf{Evaluate these definite integrals:}

\textbf{77.} \(\int_0^2 x^2 \, dx\)

\begin{itemize}
\tightlist
\item
  Antiderivative: \(F(x) = \frac{x^3}{3}\)
\item
  \(F(2) - F(0) = \frac{8}{3} - 0 ={}\) \_\_\_\_\_
\end{itemize}

\textbf{78.} \(\int_1^3 2x \, dx\)

\textbf{79.} \(\int_0^4 (3t^2 - 2t) \, dt\)

\textbf{80.} \(\int_{-1}^{1} x^3 \, dx\)

\emph{Hint: What symmetry do you notice?}

\textbf{81.} \(\int_1^4 \sqrt{x} \, dx\)

\textbf{82.} \(\int_0^3 (x^2 + 1) \, dx\)

\begin{center}\rule{0.5\linewidth}{0.5pt}\end{center}

\end{document}
